\subsubsection*{Solution}
Let

\[
\rho_{\mathrm{GHZ}} = |\mathrm{GHZ}_3\rangle\langle\mathrm{GHZ}_3|, \qquad |\mathrm{GHZ}_3\rangle = \frac{|000\rangle+|111\rangle}{\sqrt2}.
\]

This is the standard three-qubit GHZ convention used in \href{\detokenize{https://quantum.cloud.ibm.com/learning/en/courses/basics-of-quantum-information/multiple-systems/quantum-information}}{IBM Quantum's documentation}. I interpret $\psi^{\otimes4}$ as $\rho_{\mathrm{GHZ}}^{\otimes4}$, since a trace requires an operator rather than a ket. The Haar measure $dU$ is normalized.

\paragraph*{Step-by-Step Derivation}

\subparagraph*{1. Expand the projector $S$}

Using the Pauli matrix $Z$,

\[
S=|00\rangle\langle00|+|11\rangle\langle11| =\frac{1}{2}\left(I\otimes I+Z\otimes Z\right).
\]

For any $U\in U(2)$, write

\[
UZU^\dagger=\boldsymbol r\cdot\boldsymbol\sigma =r_xX+r_yY+r_zZ,
\]

where $\boldsymbol r$ is uniformly distributed on the unit sphere. Therefore,

\[
U^{\otimes4}(S\otimes S)(U^\dagger)^{\otimes4} = \frac14\left[ I^{\otimes4} +R_1R_2 +R_3R_4 +R_1R_2R_3R_4 \right],
\]

with $R=\boldsymbol r\cdot\boldsymbol\sigma$.

Higher-copy Haar twirls can generally be evaluated through Schur-Weyl duality and Weingarten calculus, as summarized in these \href{\detokenize{https://caltech-quantum-learning-theory.github.io/materials/Note_Learning_Predicting_States.pdf}}{Caltech lecture notes}. Here the Bloch-sphere moments give a simpler evaluation.

\subparagraph*{2. Evaluate the Bloch-sphere moments}

Rotational invariance gives

\[
\mathbb E[r_a r_b]=\frac{\delta_{ab}}{3},
\]

and

\[
\mathbb E[r_a r_b r_c r_d] = \frac{ \delta_{ab}\delta_{cd} +\delta_{ac}\delta_{bd} +\delta_{ad}\delta_{bc} }{15}.
\]

The factor $1/15$ follows by contracting both sides with $\delta_{ab}\delta_{cd}$ and using $\sum_a r_a^2=1$.

Consequently,

\[
\adjustbox{max width=\linewidth}{$\displaystyle {
\begin{aligned}
N=\frac14\Bigg[
&I^{\otimes4}
+\frac13\sum_{a=x,y,z}
\left(
\sigma_a\sigma_aII
+
II\sigma_a\sigma_a
\right)\\
&+\frac1{15}\sum_{a,b=x,y,z}
\left(
\sigma_a\sigma_a\sigma_b\sigma_b
+\sigma_a\sigma_b\sigma_a\sigma_b
+\sigma_a\sigma_b\sigma_b\sigma_a
\right)
\Bigg].
\end{aligned}}$}
\]

Here tensor-product symbols inside Pauli strings are suppressed.

The $28$ distinct Pauli strings in $N$ fall into the following classes:

\begingroup
\setlength{\tabcolsep}{3pt}
\begin{longtable}{@{}>{\raggedright\arraybackslash}p{0.225\linewidth} >{\raggedright\arraybackslash}p{0.225\linewidth} >{\raggedright\arraybackslash}p{0.225\linewidth} >{\raggedright\arraybackslash}p{0.225\linewidth}@{}}
\toprule
Class & Strings & Number & Coefficient \\
\midrule
\endfirsthead
\toprule
Class & Strings & Number & Coefficient \\
\midrule
\endhead
$\mathcal I$ & $IIII$ & $1$ & $1/4$ \\
$\mathcal A$ & $\sigma_a\sigma_aII$ & $3$ & $1/12$ \\
$\mathcal B$ & $II\sigma_a\sigma_a$ & $3$ & $1/12$ \\
$\mathcal D$ & $\sigma_a^{\otimes4}$ & $3$ & $1/20$ \\
$\mathcal M$ & mixed strings with two $\sigma_a$ and two $\sigma_b$, $a\neq b$ & $18$ & $1/60$ \\
\bottomrule
\end{longtable}
\endgroup

For $\mathcal D$, the same string occurs once in each of the three fourth-moment pairings, hence its coefficient is $3/60=1/20$.

\subparagraph*{3. Pauli correlations of the GHZ state}

The density operator has the stabilizer expansion

\[
\rho_{\mathrm{GHZ}} = \frac18\left( III+ZZI+ZIZ+IZZ +XXX-XYY-YXY-YYX \right).
\]

Therefore the only nonzero three-qubit Pauli expectation values are

\[
\begin{aligned}
\langle III\rangle
&=\langle ZZI\rangle
=\langle ZIZ\rangle
=\langle IZZ\rangle
=\langle XXX\rangle=1,\\
\langle XYY\rangle
&=\langle YXY\rangle
=\langle YYX\rangle=-1.
\end{aligned}
\]

Every other Pauli word has zero expectation.

\subparagraph*{4. Regroup the tensor factors by copies}

Suppose a term from row $r$ of $N$ is

\[
c_r\,P_{r1}\otimes P_{r2}\otimes P_{r3}\otimes P_{r4}.
\]

After regrouping the $12$ qubits by the four GHZ copies, its contribution is

\[
c_1c_2c_3 \prod_{j=1}^{4} \left\langle P_{1j}\otimes P_{2j}\otimes P_{3j} \right\rangle_{\mathrm{GHZ}}.
\]

Thus a row triple survives exactly when each of its four column words belongs to

\[
\mathcal G= \{III,ZZI,ZIZ,IZZ,XXX,XYY,YXY,YYX\}.
\]

Checking the $28^3$ row triples against this eight-element set leaves $112$ nonzero terms. Grouping them by the classes above gives:

\begingroup
\setlength{\tabcolsep}{3pt}
\begin{longtable}{@{}>{\raggedright\arraybackslash}p{0.300\linewidth} >{\raggedright\arraybackslash}p{0.300\linewidth} >{\raggedright\arraybackslash}p{0.300\linewidth}@{}}
\toprule
Row-class multiset & Surviving ordered triples & Total contribution \\
\midrule
\endfirsthead
\toprule
Row-class multiset & Surviving ordered triples & Total contribution \\
\midrule
\endhead
$\mathcal I,\mathcal I,\mathcal I$ & $1$ & $1/64$ \\
$\mathcal A,\mathcal A,\mathcal A$ & $4$ & $1/432$ \\
$\mathcal B,\mathcal B,\mathcal B$ & $4$ & $1/432$ \\
$\mathcal A,\mathcal A,\mathcal I$ & $3$ & $1/192$ \\
$\mathcal B,\mathcal B,\mathcal I$ & $3$ & $1/192$ \\
$\mathcal A,\mathcal B,\mathcal D$ & $6$ & $1/480$ \\
$\mathcal A,\mathcal M,\mathcal M$ & $12$ & $1/3600$ \\
$\mathcal B,\mathcal M,\mathcal M$ & $12$ & $1/3600$ \\
$\mathcal D,\mathcal D,\mathcal D$ & $4$ & $1/2000$ \\
$\mathcal D,\mathcal D,\mathcal I$ & $3$ & $3/1600$ \\
$\mathcal D,\mathcal M,\mathcal M$ & $36$ & $1/2000$ \\
$\mathcal M,\mathcal M,\mathcal M$ & $24$ & $-1/9000$ \\
\bottomrule
\end{longtable}
\endgroup

The negative final entry comes from the $XYY$, $YXY$, and $YYX$ GHZ correlations.

\subparagraph*{5. Sum the contributions}

Hence

\[
\adjustbox{max width=\linewidth}{$\displaystyle \begin{aligned}
\operatorname{tr}
\left(
N^{\otimes3}\rho_{\mathrm{GHZ}}^{\otimes4}
\right)
={}&
\frac1{64}
+2\left(
\frac1{432}
+\frac1{192}
+\frac1{3600}
\right)
+\frac1{480}\\
&+\frac1{2000}
+\frac3{1600}
+\frac1{2000}
-\frac1{9000}\\
={}&\frac{487}{13500}.
\end{aligned}$}
\]

Numerically,

\[
\frac{487}{13500} = 0.036074074074074\ldots
\]
